\documentclass[11pt,a4paper]{article}
\usepackage[T1]{fontenc}
\usepackage[utf8]{inputenc}
\usepackage[margin=23mm,headheight=15pt]{geometry}
\usepackage{amsmath,amssymb,mathtools,amsthm}
\usepackage{newpxtext,newpxmath}
\usepackage{microtype,xcolor,enumitem,environ,needspace,fancyhdr,etoolbox}
\usepackage[colorlinks=true,linkcolor=heading,urlcolor=heading]{hyperref}
\definecolor{heading}{HTML}{244B65}
\definecolor{hintcolor}{HTML}{765638}
\setlength{\parindent}{0pt}
\setlength{\parskip}{5pt plus 1pt}
\setlength{\emergencystretch}{2em}
\setlist{topsep=4pt,itemsep=4pt,parsep=2pt,leftmargin=*}
\allowdisplaybreaks[1]
\newtheoremstyle{exostyle}{8pt}{6pt}{\normalfont}{}{\color{heading}\bfseries}{.}{\newline}{}
\theoremstyle{exostyle}
\newtheorem{exercise}{Exercise}
\BeforeBeginEnvironment{exercise}{\par\addvspace{10pt}\Needspace{9\baselineskip}%
  \begingroup\color{heading!45}\hrule height .6pt\endgroup\nobreak}
\newif\ifcorr
\ifdefined\StatementOnly\corrfalse\else\corrtrue\fi
\newcommand{\showsolutions}{\corrtrue}
\newcommand{\hidesolutions}{\corrfalse}
\NewEnviron{hints}{\ifcorr\par\Needspace{4\baselineskip}\noindent
  {\color{hintcolor}\bfseries Useful hints}\par\BODY\par\medskip\fi}
\NewEnviron{solution}{\ifcorr\par\Needspace{4\baselineskip}\noindent
  {\color{heading}\bfseries Detailed solution}\par\BODY\par\medskip\fi}
\newcommand{\R}{\mathbb{R}}
\newcommand{\push}{\#}
\newcommand{\W}{\mathcal{W}}
\newcommand{\Wone}{\mathcal{W}_1}
\newcommand{\Wtwo}{\mathcal{W}_2}
\newcommand{\Unif}{\mathcal U}
\newcommand{\ip}[2]{\langle #1,#2\rangle}
\newcommand{\gradW}{\mathord{\nabla\mkern-9.3mu\nabla}}
\pagestyle{fancy}
\fancyhf{}
\fancyhead[L]{\small\color{heading}Optimal Transport / Exercises}
\fancyhead[R]{\small\ifcorr Solutions\else Statements\fi}
\fancyfoot[C]{\small\thepage}
\renewcommand{\headrulewidth}{0.3pt}
\renewcommand{\maketitle}{\thispagestyle{plain}\begingroup
  \color{heading}\LARGE\bfseries\SheetTitle\par\endgroup
  \vspace{3pt}{\large\ifcorr Statements, hints and worked solutions\else Exercise statements\fi}\par
  \vspace{7pt}\hrule\vspace{9pt}
  \ifcorr Each exercise is followed by question-by-question hints and a worked solution.
  Try the questions before reading the hints.\else
  This version contains the exercises only. Hints and worked solutions are available in the companion PDF.\fi
  \par\medskip}
\hypersetup{pdfauthor={Optimal Transport for Machine Learners},pdfsubject={Exercises in optimal transport}}

\newcommand{\SheetTitle}{Flow Matching and Dynamic OT}
\newcommand{\EE}{\mathbb{E}}
\newcommand{\Id}{\operatorname{Id}}
\newcommand{\Law}{\operatorname{Law}}
\newcommand{\Cov}{\operatorname{Cov}}
\newcommand{\Gauss}{\mathcal{N}}
\hypersetup{pdftitle={Flow Matching and Dynamic OT -- Exercises},
  pdfauthor={Julie Delon (original exercises); OT4ML (adaptation and solutions)}}
\begin{document}
\maketitle
\textbf{Original exercises: Julie Delon.} Adapted for OT4ML, with clarified
assumptions and additional hints and worked solutions.

\textbf{Setting.} All random vectors have finite second moments.
For a coupling $(X_0,X_1)$ of $\alpha$ and $\beta$ on $\R^d$, set
\[
  X_t=(1-t)X_0+tX_1,\qquad
  \alpha_t=\Law(X_t),\qquad
  v_t(x)=\EE[X_1-X_0\mid X_t=x].
\]
Conditional expectations are defined $\alpha_t$-almost everywhere. Whenever
an ODE is considered, we use a representative of $v$ satisfying the stated
regularity assumptions. For $t<1$, an equivalent formula is
\begin{equation}\label{eq-fm-conditional}
  v_t(x)=\frac{\EE[X_1\mid X_t=x]-x}{1-t}.
\end{equation}
The associated flow, when well defined, solves
$\partial_t\Phi_t(x)=v_t(\Phi_t(x))$, $\Phi_0(x)=x$.
The random straight paths $X_t$ and the ODE paths $\Phi_t(X_0)$ need not
coincide; under suitable regularity they have the same one-time laws.

\begin{exercise}[Affine changes and a time change]\label{ex-fm-affine}
Let $A\in\R^{d\times d}$ be invertible, $b\in\R^d$, and $c>0$.
Prove the following transformation rules for the conditional velocity.
\begin{enumerate}[label=(\alph*)]
\item For the coupling $(AX_0+b,AX_1+b)$, show that
  $v_t^{A,b}(x)=A v_t(A^{-1}(x-b))$.
\item For the coupling $(X_0,X_1+b)$, show that
  $v_t^b(x)=v_t(x-tb)+b$.
\item For the coupling $(X_0,cX_1)$, show that
\[
  v_t^c(x)=\frac{c}{s_t}v_{r_t}\!\left(\frac{x}{s_t}\right)
      +\frac{c-1}{s_t}x,
  \qquad s_t=1-t+tc,\qquad r_t=\frac{tc}{s_t}.
\]
\end{enumerate}
All equalities are understood almost everywhere for the corresponding
interpolated laws.
\end{exercise}
\begin{hints}
\begin{enumerate}[label=(\alph*)]
\item The new interpolant is $AX_t+b$; conditioning on it is equivalent
  to conditioning on $X_t$.
\item The new interpolant is $X_t+tb$, while its time derivative is
  $X_1-X_0+b$.
\item Write $(1-t)X_0+tcX_1=s_tX_{r_t}$ and compute $r_t'$.
\end{enumerate}
\end{hints}
\begin{solution}
\textbf{(a) Common affine change.}
Write $Y_t=AX_t+b$. Invertibility gives $X_t=A^{-1}(Y_t-b)$, so
\[
  \EE[A(X_1-X_0)\mid Y_t=x]
     =A\EE[X_1-X_0\mid X_t=A^{-1}(x-b)].
\]
This is precisely the claimed formula. It concerns flow matching,
not preservation of quadratic optimality; Exercise~\ref{ex-fm-optimality}
explains the distinction.

\textbf{(b) Translation of the target.}
Now $Y_t=X_t+tb$ and $\dot Y_t=X_1-X_0+b$. Consequently,
\[
  v_t^b(x)=\EE[X_1-X_0+b\mid X_t=x-tb]=v_t(x-tb)+b.
\]

\textbf{(c) Dilation of the target.}
Since $s_t>0$, the identity $Y_t=s_tX_{r_t}$ is valid throughout $[0,1]$.
Moreover, $s_t'=c-1$ and $r_t'=c/s_t^2$. Differentiating the identity
for each realization of $(X_0,X_1)$ yields
\[
  \dot Y_t=(c-1)X_{r_t}+\frac{c}{s_t}(X_1-X_0).
\]
Conditioning on $Y_t=x$, equivalently on $X_{r_t}=x/s_t$, proves the
formula. When $c=1$, both the dilation and the time change disappear.
\end{solution}

\begin{exercise}[Which transformations preserve optimality?]\label{ex-fm-optimality}
Assume that $\alpha$ has a density, and let $T^*$ be its optimal transport
map to $\beta$ for the cost $|x-y|^2$. Apply the transformations of
Exercise~\ref{ex-fm-affine} to the optimal coupling $(X_0,T^*(X_0))$.
\begin{enumerate}[label=(\alph*)]
\item For the common affine transformation $S(x)=Ax+b$, identify the
  induced map from $S_\push\alpha$ to $S_\push\beta$. Prove that it is
  optimal when $A=sO$, with $s>0$ and $O$ orthogonal. Is this true for
  every invertible $A$? If not, give a two-dimensional counterexample.
\item Identify the optimal map from $\alpha$ to $(x\mapsto x+b)_\push\beta$.
\item Identify the optimal map from $\alpha$ to $(x\mapsto cx)_\push\beta$
  for $c>0$.
\end{enumerate}
\end{exercise}
\begin{hints}
\begin{enumerate}[label=(\alph*)]
\item The candidate is $S\circ T^*\circ S^{-1}$. Similarities multiply
  every squared distance by $s^2$. For a counterexample, conjugate a
  symmetric positive definite linear map by a shear.
\item Write $T^*=\nabla\varphi$ with $\varphi$ convex and add a linear
  function to $\varphi$.
\item Multiply the same potential by $c$.
\end{enumerate}
\end{hints}
\begin{solution}
\textbf{(a) Common change of coordinates.}
The induced map is
\[
  \widetilde T(x)=A T^*(A^{-1}(x-b))+b.
\]
It has the required pushforward property for every invertible $A$.
If $A=sO$, then $|S(x)-S(y)|^2=s^2|x-y|^2$. Applying $S$ to both
coordinates gives a bijection between the two sets of couplings and
multiplies all their costs by the same constant. Optimality is therefore
preserved.

For a general $A$, the transformed cost is instead $|A(x-y)|^2$;
this changes the geometry. For an explicit counterexample, take
\[
  \alpha=\Gauss(0,I_2),\qquad T^*(x)=Bx,\qquad
  B=\begin{pmatrix}2&0\\0&1\end{pmatrix},\qquad
  A=\begin{pmatrix}1&1\\0&1\end{pmatrix},\qquad b=0,
\]
and set $\beta=T^*_\push\alpha$. The original map is optimal because
$B$ is symmetric positive definite. The conjugated map has matrix
\[
  ABA^{-1}=\begin{pmatrix}2&-1\\0&1\end{pmatrix}.
\]
This matrix is not symmetric, so the map cannot be the gradient of a
convex function. Since the transformed source has a strictly positive
density, Brenier's theorem excludes its optimality. In dimension one
this obstruction does not occur: every invertible scalar is a similarity.

\textbf{(b) Target translation.}
The map $T^*+b$ has the correct target law and is the gradient of the
convex function $\varphi(x)+b\cdot x$. It is therefore the optimal map.
One can also expand the cost: the change
$|x-y-b|^2-|x-y|^2=-2b\cdot(x-y)+|b|^2$ has an integral depending
only on the marginals.

\textbf{(c) Positive target dilation.}
The map $cT^*$ pushes $\alpha$ to the dilated target and equals
$\nabla(c\varphi)$. Positivity of $c$ preserves convexity, so this map
is optimal. Equivalently,
\[
  |x-cy|^2=c|x-y|^2+(1-c)|x|^2+(c^2-c)|y|^2;
\]
the last two terms have fixed integrals and the factor $c$ is positive.
\end{solution}

\begin{exercise}[One-dimensional flow matching gives an optimal terminal map]\label{ex-fm-line}
Let $\alpha,\beta$ be probability measures on $\R$ with finite second
moments, and let $(X_0,X_1)$ be any coupling of them.
Assume that $v$ has a continuous representative on $[0,1)\times\R$ such
that, for each $T<1$, there is $L_T<\infty$ with
\[
  |v_t(x)-v_t(y)|\leq L_T|x-y|,\qquad |v_t(0)|\leq L_T
  \quad (0\leq t\leq T).
\]
Let $\Phi_t$ denote the ODE flow for $t<1$. No bound on $L_T$ as
$T\uparrow1$ is assumed.
\begin{enumerate}[label=(\alph*)]
\item Show that $\alpha_t$ solves
  $\partial_t\alpha_t+\partial_x(v_t\alpha_t)=0$ in the weak sense,
  and deduce $(\Phi_t)_\push\alpha=\alpha_t$ for $t<1$.
\item Show that $\Phi_t$ is strictly increasing for $t<1$ and that
  $\Phi_1(x):=\lim_{t\uparrow1}\Phi_t(x)$ exists for $\alpha$-almost
  every $x$. Can two distinct trajectories have the same endpoint?
\item Prove that $\Phi_1$ pushes $\alpha$ to $\beta$ and is optimal for
  the squared-distance cost. When $\alpha$ has no atoms, identify it
  using cumulative and quantile functions.
\end{enumerate}
\end{exercise}
\begin{hints}
\begin{enumerate}[label=(\alph*)]
\item Differentiate $\EE[\psi(X_t)]$ for a smooth compactly supported
  test function $\psi$. Use uniqueness of the continuity equation
  associated with a Lipschitz velocity field.
\item Use backward uniqueness to rule out a meeting before time $1$.
  For existence of the terminal limit, apply conditional Jensen to
  $v_t(X_t)$ and bound the length of almost every characteristic.
\item Compare $X_t$ with $X_1$. Use bounded continuous test functions
  to pass to the limit in $(\Phi_t)_\push\alpha=\alpha_t$, then use
  optimality of monotone one-dimensional couplings.
\end{enumerate}
\end{hints}
\begin{solution}
\textbf{(a) The prescribed marginals.}
Differentiation under the expectation is justified by boundedness of
$\psi'$ and integrability of $|X_1-X_0|$. Thus
\[
  \frac{d}{dt}\int\psi\,d\alpha_t
    =\EE[\psi'(X_t)(X_1-X_0)]
    =\int\psi'(x)v_t(x)\,d\alpha_t(x).
\]
This is the weak continuity equation. On every interval $[0,T]$, the
assumptions give a unique global characteristic flow and uniqueness
of the corresponding narrowly continuous measure solution. The
characteristic representation therefore gives
$(\Phi_t)_\push\alpha=\alpha_t$. One way to see uniqueness is to test
against functions transported backward by the flow: their integral
against any solution is constant. For Lipschitz velocities this
argument follows by approximation from the smooth case.

\textbf{(b) Noncrossing characteristics.}
Fix $x<y$. If $\Phi_t(x)$ and $\Phi_t(y)$ first met before time $1$,
uniqueness backward from their meeting point would imply $x=y$.
Continuity of the trajectories then gives $\Phi_t(x)<\Phi_t(y)$ for
all $t<1$. To reach time $1$, we need a bound on path length, not a
uniform Lipschitz constant. Part~(a), Tonelli's theorem and conditional
Jensen give
\[
 \int\!\int_0^1 |v_t(\Phi_t(x))|^2\,dt\,d\alpha(x)
 =\int_0^1\EE\bigl[|\EE[X_1-X_0\mid X_t]|^2\bigr]\,dt
 \le \EE|X_1-X_0|^2<\infty .
\]
For almost every initial point the speed is therefore square integrable,
hence integrable on $(0,1)$. Its tail integral bounds the displacement
between any two late times and tends to zero. The trajectory is Cauchy
and has a finite limit. On this full-measure set the limit map is
nondecreasing, but it need not be strictly increasing: for $X_1=0$
the trajectories $\Phi_t(x)=(1-t)x$ all end at the origin.

\textbf{(c) Terminal transport and optimality.}
First we show that the prescribed intermediate law approaches the target.
Let $\pi=\Law(X_0,X_1)$ and consider the new coupling
\[
  \pi_t=\bigl((x_0,x_1)\mapsto((1-t)x_0+tx_1,x_1)\bigr)_\push\pi
       =\Law(X_t,X_1).
\]
Its first marginal is $\alpha_t$ by definition of $X_t$, and its second
is $\beta$ because $X_1\sim\beta$. It is therefore an admissible
competitor in the definition of $\Wtwo(\alpha_t,\beta)$, whether or
not the original coupling $\pi$ is optimal. Consequently,
\begin{align*}
  \Wtwo(\alpha_t,\beta)^2
    &=\inf_{\gamma\text{ coupling }\alpha_t,\beta}
                       \int|x-y|^2\,d\gamma(x,y)\\
    &\leq\int|x-y|^2\,d\pi_t(x,y)
     =\EE|X_t-X_1|^2\\
    &=\EE|(1-t)X_0+tX_1-X_1|^2
     =(1-t)^2\EE|X_0-X_1|^2.
\end{align*}
The last expectation is finite, since
$\EE|X_0-X_1|^2\leq2\EE|X_0|^2+2\EE|X_1|^2$.
Thus $\Wtwo(\alpha_t,\beta)\to0$, and in particular
$\alpha_t$ converges weakly to $\beta$.

We must next identify the law of the terminal \emph{ODE} map.
Choose a bounded continuous test function $\psi$. By the
almost-everywhere convergence $\Phi_t(x)\to\Phi_1(x)$ and the bound
$|\psi(\Phi_t(x))|\leq\|\psi\|_\infty$, dominated convergence gives
\begin{align*}
  \int\psi(\Phi_1(x))\,d\alpha(x)
    &=\lim_{t\uparrow1}\int\psi(\Phi_t(x))\,d\alpha(x)\\
    &=\lim_{t\uparrow1}\int\psi(y)\,d\alpha_t(y)
     =\int\psi(y)\,d\beta(y).
\end{align*}
The middle equality is the pushforward identity proved in part~(a);
the last is weak convergence. Since this holds for every bounded
continuous $\psi$, $(\Phi_1)_\push\alpha=\beta$.

Finally, the limiting map is nondecreasing by part~(b), so its coupling
is the monotone, or quantile, coupling. Write
$\mathcal C_\eta(x)=\eta(({-\infty},x])$ and
$\mathcal C_\eta^{-1}(r)=\inf\{x:\mathcal C_\eta(x)\geq r\}$ for
$0<r<1$ and $\eta=\alpha,\beta$.
To see the identification directly, for $X\sim\alpha$ the events $\{X\leq a\}$ and
$\{\Phi_1(X)\leq b\}$ are nested initial segments of the ordered
full-measure set on which $\Phi_1$ is defined. Hence
\[
  \Pr\bigl(X\leq a,\ \Phi_1(X)\leq b\bigr)
       =\min\{\mathcal C_\alpha(a),\mathcal C_\beta(b)\}.
\]
This is also the joint cumulative function of
$(\mathcal C_\alpha^{-1}(U),\mathcal C_\beta^{-1}(U))$ for
$U$ uniform on $(0,1)$. The one-dimensional rearrangement theorem
therefore gives
\[
  \int|x-\Phi_1(x)|^2\,d\alpha(x)
    =\int_0^1|\mathcal C_\alpha^{-1}(r)-\mathcal C_\beta^{-1}(r)|^2\,dr
    =\Wtwo(\alpha,\beta)^2.
\]
In particular, when $\alpha$ is atomless,
\[
  \Phi_1(x)=\mathcal C_\beta^{-1}(\mathcal C_\alpha(x))
  \quad\text{for $\alpha$-almost every $x$},
\]
because $\mathcal C_\alpha(X)$ is then uniform on $(0,1)$.
This conclusion concerns the \emph{terminal map}, not the
original coupling or optimality of the entire path. With atoms in
$\alpha$, the same monotonicity argument applies whenever the flow
assumptions hold, but the displayed quantile composition need not be
a valid formula for arbitrary atomic source and target laws: a
deterministic flow cannot split a source atom.
\end{solution}

\begin{exercise}[Jointly Gaussian endpoints]\label{ex-fm-gaussian}
Suppose that the endpoints are jointly centered Gaussian, with
\[
  (X_0,X_1)\sim\Gauss\!\left(0,
     \begin{pmatrix}\Sigma_0&R^\top\\R&\Sigma_1\end{pmatrix}\right),
  \qquad R=\EE[X_1X_0^\top],\qquad \Sigma_0,\Sigma_1\succ0.
\]
The joint covariance is positive semidefinite. In addition, assume
that $\Sigma_t:=\Cov(X_t)$ is positive definite for every $t\in[0,1]$.
A sufficient condition is positive definiteness of the joint covariance.
For jointly Gaussian vectors $(Z,Y)$ with $\Cov(Y)$ invertible,
recall the Gaussian regression formula
\[
  \EE[Z\mid Y=y]=\EE Z+\Cov(Z,Y)\Cov(Y)^{-1}(y-\EE Y).
\]
\begin{enumerate}[label=(\alph*)]
\item Compute $\Sigma_t$ and prove that $v_t(x)=B_tx$, where
\begin{equation}\label{eq-fm-gaussian-velocity}
  B_t=\big[(1-t)(R-\Sigma_0)+t(\Sigma_1-R^\top)\big]\Sigma_t^{-1}.
\end{equation}
  Recover, for $t<1$, the equivalent expression
  $B_t=\big([(1-t)R+t\Sigma_1]\Sigma_t^{-1}-I_d\big)/(1-t)$.
\item Assume now that $X_0$ and $X_1$ are independent and that
  $\Sigma_0\Sigma_1=\Sigma_1\Sigma_0$. Solve the ODE explicitly and
  show that its terminal map is the quadratic optimal map between
  $\Gauss(0,\Sigma_0)$ and $\Gauss(0,\Sigma_1)$.
  Is the whole curve $\alpha_t$ necessarily a Wasserstein geodesic?
\end{enumerate}
\end{exercise}
\begin{hints}
\begin{enumerate}[label=(\alph*)]
\item Expand $\Cov(X_t)$ and $\Cov(X_1-X_0,X_t)$, keeping track of
  the order of the cross-covariance matrices.
\item Independence gives $R=0$. Simultaneously diagonalize the two
  covariances and solve each scalar ODE. For the last question, take
  $\Sigma_0=\Sigma_1=I_d$.
\end{enumerate}
\end{hints}
\begin{solution}
\textbf{Gaussian regression: proof of the recalled formula.}
For column vectors, our convention is
$\Cov(Z,Y)=\EE[(Z-\EE Z)(Y-\EE Y)^\top]$.
Write $m_Z=\EE Z$, $m_Y=\EE Y$, $Q=\Cov(Y)$ and
$L=\Cov(Z,Y)Q^{-1}$. Subtract the proposed affine predictor from $Z$
and define the centered residual
\[
  \eta=Z-m_Z-L(Y-m_Y).
\]
The pair $(\eta,Y-m_Y)$ is jointly Gaussian because it is a linear
transformation of the centered jointly Gaussian pair $(Z-m_Z,Y-m_Y)$.
Its cross covariance vanishes:
\begin{align*}
  \Cov(\eta,Y)
    &=\EE\bigl[(Z-m_Z-L(Y-m_Y))(Y-m_Y)^\top\bigr]\\
    &=\Cov(Z,Y)-LQ\\
    &=\Cov(Z,Y)-\Cov(Z,Y)Q^{-1}Q=0.
\end{align*}
For jointly Gaussian vectors, zero cross covariance implies
independence. Indeed, the joint characteristic function factors:
\[
  \EE e^{\,i u^\top\eta+i v^\top(Y-m_Y)}
   =e^{-\frac12u^\top\Cov(\eta)u-\frac12v^\top Qv}
   =\EE e^{\,i u^\top\eta}\,\EE e^{\,i v^\top(Y-m_Y)}.
\]
This argument also allows a singular residual covariance.
Since $\eta$ is independent of $Y$ and $\EE\eta=0$, conditioning the
identity $Z=m_Z+L(Y-m_Y)+\eta$ on $Y$ gives
\[
  \EE[Z\mid Y]
     =m_Z+L(Y-m_Y)+\EE[\eta\mid Y]
     =m_Z+\Cov(Z,Y)\Cov(Y)^{-1}(Y-m_Y).
\]
Evaluating this affine version at $Y=y$ proves the recalled formula.
Only $\Cov(Y)$ must be invertible; the joint covariance need not be.

\par\Needspace{5\baselineskip}
\textbf{(a) Regression and covariance.}
All the vectors in this question are centered, so covariance is simply
the expectation of an outer product. In particular,
\[
  \EE[X_0X_0^\top]=\Sigma_0,\quad
  \EE[X_1X_1^\top]=\Sigma_1,\quad
  \EE[X_1X_0^\top]=R,\quad
  \EE[X_0X_1^\top]=R^\top.
\]
Substitute $X_t=(1-t)X_0+tX_1$ into $\Sigma_t=\EE[X_tX_t^\top]$
and distribute the product:
\begin{align*}
  \Sigma_t
    &=\EE\bigl[((1-t)X_0+tX_1)((1-t)X_0+tX_1)^\top\bigr]\\
    &=(1-t)^2\EE[X_0X_0^\top]
       +t(1-t)\EE[X_0X_1^\top]\\
    &\hspace{1.5em}+t(1-t)\EE[X_1X_0^\top]+t^2\EE[X_1X_1^\top].
\end{align*}
Replacing these four expectations by their matrix values yields
\begin{equation}\label{eq-fm-gaussian-covariance}
  \Sigma_t=(1-t)^2\Sigma_0+t(1-t)(R+R^\top)+t^2\Sigma_1.
\end{equation}
For the covariance needed in the regression, the same expansion gives
\begin{align*}
  \Cov(X_1-X_0,X_t)
    &=\EE\bigl[(X_1-X_0)((1-t)X_0+tX_1)^\top\bigr]\\
    &=(1-t)\EE[(X_1-X_0)X_0^\top]
         +t\EE[(X_1-X_0)X_1^\top]\\
    &=(1-t)\bigl(\EE[X_1X_0^\top]-\EE[X_0X_0^\top]\bigr)\\
    &\hspace{1.5em}+t\bigl(\EE[X_1X_1^\top]-\EE[X_0X_1^\top]\bigr)\\
    &=(1-t)(R-\Sigma_0)+t(\Sigma_1-R^\top).
\end{align*}
The transpose in the last term is essential: the order of the two
vectors in a cross covariance cannot be interchanged without
transposing the matrix.

Now apply the regression formula with $Z=X_1-X_0$ and $Y=X_t$.
These two vectors are jointly Gaussian, their means are zero,
and $\Cov(Y)=\Sigma_t$ is invertible by assumption. Therefore
\begin{align*}
  v_t(x)
    &=\EE[X_1-X_0\mid X_t=x]\\
    &=\Cov(X_1-X_0,X_t)\Sigma_t^{-1}x\\
    &=\big[(1-t)(R-\Sigma_0)+t(\Sigma_1-R^\top)\big]\Sigma_t^{-1}x
     =B_tx.
\end{align*}
This proves~\eqref{eq-fm-gaussian-velocity}, including the order
of the matrix factors.

For the alternative expression, compute instead
\[
  \Cov(X_1,X_t)
     =\EE[X_1((1-t)X_0+tX_1)^\top]
     =(1-t)R+t\Sigma_1.
\]
Regression gives
$\EE[X_1\mid X_t=x]=[(1-t)R+t\Sigma_1]\Sigma_t^{-1}x$.
Substituting in~\eqref{eq-fm-conditional}, valid for $t<1$, gives
\[
  v_t(x)
    =\frac{[(1-t)R+t\Sigma_1]\Sigma_t^{-1}x-x}{1-t}
    =\frac{[(1-t)R+t\Sigma_1]\Sigma_t^{-1}-I_d}{1-t}\,x.
\]
Formula~\eqref{eq-fm-gaussian-velocity} avoids a removable denominator
at $t=1$.

The extra invertibility assumption matters. Positive definite
\emph{endpoint} covariances alone do not imply it: if $X_1=-X_0$
with $X_0\sim\Gauss(0,I_d)$, then $X_{1/2}=0$ and
$\Sigma_{1/2}=0$. A conditional velocity still exists almost everywhere
in this example, but the inverse-covariance formula is not valid there.

\textbf{(b) Independent, commuting covariances.}
Now $R=0$, so
\[
  \Sigma_t=(1-t)^2\Sigma_0+t^2\Sigma_1,\qquad
  B_t=\big[t\Sigma_1-(1-t)\Sigma_0\big]\Sigma_t^{-1}.
\]
Choose an orthogonal matrix $O$ with
$\Sigma_i=O\operatorname{diag}(\lambda_{i,1},\ldots,\lambda_{i,d})O^\top$.
All eigenvalues $\lambda_{i,j}$ are positive. In these coordinates,
$z(t)=O^\top\Phi_t(x)$ has initial value $z(0)=O^\top x$, and the
$j$th component of its ODE is
\[
  \dot z_j(t)=
  \frac{t\lambda_{1,j}-(1-t)\lambda_{0,j}}
       {(1-t)^2\lambda_{0,j}+t^2\lambda_{1,j}}\,z_j(t).
\]
Set $a_j(t)=(1-t)^2\lambda_{0,j}+t^2\lambda_{1,j}>0$. Its derivative is
$a_j'(t)=2[t\lambda_{1,j}-(1-t)\lambda_{0,j}]$, so the equation reads
$\dot z_j(t)=a_j'(t)z_j(t)/(2a_j(t))$.
An integrating factor is $a_j(t)^{-1/2}$, because
\[
  \frac{d}{dt}\left(\frac{z_j(t)}{\sqrt{a_j(t)}}\right)
    =\frac{\dot z_j(t)}{\sqrt{a_j(t)}}
       -\frac{a_j'(t)z_j(t)}{2a_j(t)^{3/2}}=0.
\]
Thus $z_j(t)/\sqrt{a_j(t)}$ is constant and equals
$z_j(0)/\sqrt{a_j(0)}=z_j(0)/\sqrt{\lambda_{0,j}}$.
Equivalently, the explicit integration is
\begin{align*}
  z_j(t)
    &=z_j(0)\exp\left(\int_0^t\frac{a_j'(s)}{2a_j(s)}\,ds\right)\\
    &=z_j(0)\exp\left(\frac12\log\frac{a_j(t)}{a_j(0)}\right)\\
    &=z_j(0)\sqrt{\frac{(1-t)^2\lambda_{0,j}+t^2\lambda_{1,j}}
                             {\lambda_{0,j}}}.
\end{align*}
The integrating-factor argument never divides by $z_j$ and therefore
also covers $z_j(0)=0$, in which case the solution stays zero.
Returning to the original coordinates gives
\[
  \Phi_t(x)
    =O\operatorname{diag}\left(\sqrt{\frac{a_1(t)}{\lambda_{0,1}}},
                    \ldots,\sqrt{\frac{a_d(t)}{\lambda_{0,d}}}\right)O^\top x
    =\Sigma_t^{1/2}\Sigma_0^{-1/2}x.
\]
At time $1$, the matrix $M=\Sigma_1^{1/2}\Sigma_0^{-1/2}$ is symmetric
positive definite because the covariances commute. Moreover,
$M\Sigma_0M^\top=\Sigma_1$. Thus $\Phi_1$ has the correct Gaussian
target and is the gradient of the convex potential $x^\top Mx/2$;
it is the Brenier map.

The path need not be a geodesic. For $\Sigma_0=\Sigma_1=I_d$, the
terminal map is the identity, but
$\alpha_t=\Gauss(0,[(1-t)^2+t^2]I_d)$ first contracts and then expands.
The Wasserstein geodesic between two identical laws is constant.
\end{solution}

\newpage
\begin{exercise}[Gaussian mixtures and a discrete target]\label{ex-fm-mixtures}
Let the joint law of the endpoints be a finite Gaussian mixture:
\[
  \Law(X_0,X_1)=\sum_{k=1}^K\lambda_k
   \Gauss\!\left(
     \begin{pmatrix}m_0^k\\m_1^k\end{pmatrix},
     \begin{pmatrix}\Sigma_0^k&R_k^\top\\R_k&\Sigma_1^k\end{pmatrix}
   \right),
  \qquad \lambda_k>0,\quad\sum_k\lambda_k=1.
\]
The component covariance matrices may be singular, but assume that each
interpolated covariance $\Sigma_t^k$ is positive definite for $0\leq t<1$.
\begin{enumerate}[label=(\alph*)]
\item Compute the mean $m_t^k$ and covariance $\Sigma_t^k$ of the
  interpolant in component $k$, and its conditional velocity $w_t^k$.
\item Let $p_t^k$ be the density of $\Gauss(m_t^k,\Sigma_t^k)$.
  Prove the mixture formula
\[
  v_t(x)=\sum_{k=1}^K\omega_{k,t}(x)w_t^k(x),\qquad
  \omega_{k,t}(x)=
     \frac{\lambda_kp_t^k(x)}{\sum_{j=1}^K\lambda_jp_t^j(x)}.
\]
\item Specialize to \emph{independent} endpoints
  $X_0\sim\Gauss(0,I_d)$ and
  $X_1\sim K^{-1}\sum_{k=1}^K\delta_{y_k}$.
  Give explicit formulas for $\alpha_t$, $\omega_{k,t}(x)$ and $v_t(x)$
  for $t<1$. Explain what happens to $\alpha_t$ as $t\uparrow1$.
\end{enumerate}
\end{exercise}
\begin{hints}
\begin{enumerate}[label=(\alph*)]
\item Introduce the component label $J$, with $\Pr(J=k)=\lambda_k$,
  and apply Exercise~\ref{ex-fm-gaussian} conditionally on $J=k$.
\item Apply Bayes' rule to $\Pr(J=k\mid X_t=x)$, then condition
  first on both $J$ and $X_t$.
\item In component $k$, $X_1=y_k$ is deterministic, so
  $X_t\mid(J=k)\sim\Gauss(ty_k,(1-t)^2I_d)$. Use
  \eqref{eq-fm-conditional} and cancel the common Gaussian prefactor.
\end{enumerate}
\end{hints}
\begin{solution}
\textbf{(a) Each Gaussian component.}
Conditionally on $J=k$, the interpolant has parameters
\[
  m_t^k=(1-t)m_0^k+tm_1^k,\qquad
  \Sigma_t^k=(1-t)^2\Sigma_0^k+t(1-t)(R_k+R_k^\top)+t^2\Sigma_1^k.
\]
Gaussian regression, which only needs invertibility of $\Sigma_t^k$,
gives
\[
  w_t^k(x)=m_1^k-m_0^k+B_t^k(x-m_t^k),\qquad
  B_t^k=\big[(1-t)(R_k-\Sigma_0^k)
                    +t(\Sigma_1^k-R_k^\top)\big](\Sigma_t^k)^{-1}.
\]
Equivalently, this follows from the centered formula and the
time-dependent translation of the mean.

\textbf{(b) Posterior weights.}
The density of $X_t$ is $\sum_j\lambda_jp_t^j(x)$, which is positive.
Bayes' rule therefore gives
$\Pr(J=k\mid X_t=x)=\omega_{k,t}(x)$, and the tower property yields
\[
  \EE[X_1-X_0\mid X_t=x]
  =\sum_k\Pr(J=k\mid X_t=x)
          \EE[X_1-X_0\mid X_t=x,J=k]
  =\sum_k\omega_{k,t}(x)w_t^k(x).
\]
These are \emph{posterior} weights: they generally differ from the
original mixture weights and depend on both $x$ and $t$.

\textbf{(c) Independent Gaussian noise and a discrete target.}
The component parameters are $m_0^k=0$, $m_1^k=y_k$,
$\Sigma_0^k=I_d$, $\Sigma_1^k=0$ and $R_k=0$. Thus
\[
  \alpha_t=\frac1K\sum_{k=1}^K\Gauss(ty_k,(1-t)^2I_d),\qquad
  w_t^k(x)=\frac{y_k-x}{1-t}.
\]
The Gaussian normalization factors are the same in every component,
so the responsibilities and velocity are
\[
  \omega_{k,t}(x)=
    \frac{\exp\!\left(-\dfrac{|x-ty_k|^2}{2(1-t)^2}\right)}
         {\sum_j\exp\!\left(-\dfrac{|x-ty_j|^2}{2(1-t)^2}\right)},
  \qquad
  v_t(x)=\frac{\sum_k\omega_{k,t}(x)y_k-x}{1-t}.
\]
At time $0$ the weights are all $1/K$ and
$v_0(x)=K^{-1}\sum_k y_k-x$. For any $t<1$, every component has
positive variance; its variance tends to zero and its mean tends
to $y_k$ as $t\uparrow1$. More precisely, the original coupling gives
\[
  \Wtwo(\alpha_t,\beta)^2
   \leq (1-t)^2\EE|X_0-X_1|^2
   =(1-t)^2\left(d+\frac1K\sum_k|y_k|^2\right),
  \qquad \beta=\frac1K\sum_k\delta_{y_k}.
\]
This proves convergence to the discrete target. The velocity formula
is only stated for $t<1$: the endpoint singularity allows a flow with
smooth intermediate densities to collapse onto atoms. Specifying the
two marginals alone would not determine these formulas; independence
of $X_0$ and $X_1$ is essential to this particular construction.
\end{solution}

\par\addvspace{8pt}\Needspace{5\baselineskip}
\textbf{Further reading.}
See OT4ML, Chapter~16, \emph{Generative Models via Transportation},
especially the conditional-expectation construction of flow matching
and the proposition \emph{Gaussian flow matching and optimality}.
For invariance properties and explicit Gaussian and mixture formulas,
see J.~Hertrich, A.~Chambolle and J.~Delon,
\href{https://arxiv.org/abs/2505.19712}{\emph{On the Relation between
Rectified Flows and Optimal Transport}}, NeurIPS 2025.
\end{document}
